\chapter{Geometric Constructions}\label{ch:g5:geometry}

Ruler, set square, compass: with these three tools and a bit of
method, any figure described in words can be built exactly. This
chapter constructs perpendiculars, parallels, triangles and the
special quadrilaterals --- craft first, theory later
(\cref{ch:g6:lines,ch:g6:shapes}).

\section{Perpendiculars and parallels through a point}

\begin{method}[The two basic constructions]\label{met:g5:geometry:basic}
\emph{Perpendicular to $d$ through $A$}: slide the set square along
$d$ until its other right-angle edge reaches $A$; draw; mark the
right angle.

\emph{Parallel to $d$ through $A$}: press a ruler against the set
square as a rail laid along $d$; slide the set square up to $A$;
draw. (Check: the gap between the two lines is the same
everywhere.)
\end{method}

\begin{example}[Chaining constructions]\label{ex:g5:geometry:chain}
Many figures are chains of these two moves. A rectangle $ABCD$ with
$AB = 5$ cm and $BC = 3$ cm: draw $[AB]$; perpendicular at $A$;
perpendicular at $B$; mark $D$ and $C$ at $3$ cm on the same side;
join. Every step is one basic construction.
\end{example}

\section{Triangles with the compass}

\begin{method}[Triangle from three sides]\label{met:g5:geometry:triangle}
To build a triangle with sides $6$ cm, $5$ cm, $4$ cm:
\begin{enumerate}
  \item draw the longest side, $[AB]$ with $AB = 6$ cm;
  \item arc of center $A$ and radius $5$ cm;
  \item arc of center $B$ and radius $4$ cm;
  \item their crossing is the third vertex $C$; join.
\end{enumerate}
The compass guarantees the lengths: every point of the first arc is
at exactly $5$ cm from $A$.
\end{method}

\begin{omfigure}
\begin{tikzpicture}[scale=0.95]
  \coordinate (A) at (0,0);
  \coordinate (B) at (6*0.72,0);
  \coordinate (C) at (2.7,2.376);
  \draw[thick] (A) -- (B);
  \draw[thick, omDef] (A) -- (C);
  \draw[thick, omThm] (B) -- (C);
  \draw[omDef, thin] (C) arc[start angle=41.3, delta angle=20,
    radius=3.6];
  \draw[omDef, thin] (C) arc[start angle=41.3, delta angle=-20,
    radius=3.6];
  \draw[omThm, thin] (C) arc[start angle=124.4, delta angle=20,
    radius=2.88];
  \draw[omThm, thin] (C) arc[start angle=124.4, delta angle=-20,
    radius=2.88];
  \fill (A) circle (1.5pt) node[below left, font=\small] {$A$};
  \fill (B) circle (1.5pt) node[below right, font=\small] {$B$};
  \fill (C) circle (1.5pt) node[above, font=\small] {$C$};
  \node[below, font=\small] at (2.16,0) {$6$ cm};
  \node[omDef, font=\small] at (0.85,1.45) {$5$ cm};
  \node[omThm, font=\small] at (3.9,1.35) {$4$ cm};
\end{tikzpicture}

{\small Two arcs, one crossing: the triangle $6$--$5$--$4$ built with
the compass. (If the two short sides together were shorter than the
base, the arcs would not meet --- more on this in
\cref{ch:g7:triangles}.)}
\end{omfigure}

\begin{example}[Isosceles and equilateral for free]\label{ex:g5:geometry:special}
Same compass opening for both arcs $\Rightarrow$ isosceles triangle.
Opening equal to the base $\Rightarrow$ equilateral. The compass
does not measure: it \emph{copies} lengths, which is even better.
\end{example}

\section{The special quadrilaterals, by construction}

\begin{method}[Building the family]\label{met:g5:geometry:quadrilaterals}
\begin{enumerate}
  \item \emph{Rectangle}: two pairs of parallels crossing at right
        angles --- or the chain of \cref{ex:g5:geometry:chain};
  \item \emph{square}: a rectangle whose four sides use the same
        length --- carry it with the compass;
  \item \emph{rhombus}: four equal sides --- draw one side, then two
        compass arcs from its ends with the same opening fix the
        other two vertices (a rhombus is two isosceles triangles
        glued along a diagonal).
\end{enumerate}
After building, always \emph{check} with the tools: set square on
the corners, compass on the sides.
\end{method}

\begin{example}[Recognizing from a description]\label{ex:g5:geometry:recognize}
``A quadrilateral with four $4$ cm sides and one right angle'' ---
build it: the four equal sides force a rhombus shape, and pressing
one corner to a right angle straightens all the others: the result
is a square. Some descriptions secretly describe a more special
shape than they announce. (Why the right angle spreads to all four
corners is explained in \cref{ch:g7:central}.)
\end{example}

\section{Circles in constructions}

\begin{example}[The circle as a tool]\label{ex:g5:geometry:circle}
``Find all points at $3$ cm from $A$'': the circle of center $A$,
radius $3$ cm. ``Find all points at $3$ cm from $A$ \emph{and}
$4$ cm from $B$'': draw both circles; their crossings (two, one or
zero points) are the answers. Treasure maps and triangle
constructions use exactly this idea.
\end{example}

\section{Exercises}

\begin{exercise}[$\star$]\label{exo:g5:geometry:1}
Draw a line $d$ and a point $A$ at about $3$ cm from it. Construct
the perpendicular to $d$ through $A$, then the parallel to $d$
through $A$. Mark the right angles.
\end{exercise}

\begin{exercise}[$\star$]\label{exo:g5:geometry:2}
Construct a rectangle $ABCD$ with $AB = 7$ cm and $BC = 4$ cm,
following \cref{ex:g5:geometry:chain}. Check the fourth angle with
the set square: if the construction is exact, it is right
automatically.
\end{exercise}

\begin{exercise}[$\star$]\label{exo:g5:geometry:3}
Construct a triangle with sides $7$ cm, $6$ cm and $3$ cm
(\cref{met:g5:geometry:triangle}).
\end{exercise}

\begin{exercise}[$\star$]\label{exo:g5:geometry:4}
Construct an isosceles triangle with base $5$ cm and equal sides of
$7$ cm; then an equilateral triangle of side $6$ cm.
\end{exercise}

\begin{exercise}[$\star$]\label{exo:g5:geometry:5}
Construct a square of side $5$ cm. Which tools guarantee the right
angles? Which guarantee the equal sides?
\end{exercise}

\begin{exercise}[$\star$]\label{exo:g5:geometry:6}
Construct a rhombus of side $4$ cm that is not a square (choose your
own opening for the first corner). Check the four sides with the
compass.
\end{exercise}

\begin{exercise}[$\star$]\label{exo:g5:geometry:7}
Draw two points $A$ and $B$ with $AB = 5$ cm. Construct all points
that are at $4$ cm from $A$ and at $3$ cm from $B$. How many are
there?
\end{exercise}

\begin{exercise}[$\star$]\label{exo:g5:geometry:8}
Write a construction program (numbered steps, named points) for a
rectangle $6$ cm $\times$ $2$ cm topped by an isosceles triangle
with equal sides $4$ cm (a pencil shape). Have a classmate execute
it.
\end{exercise}

\begin{exercise}[$\star\star$]\label{exo:g5:geometry:9}
Try to construct a triangle with sides $8$ cm, $3$ cm and $4$ cm.
What happens with the arcs? Explain in one sentence why this
triangle cannot exist.
\end{exercise}

\begin{exercise}[$\star\star$]\label{exo:g5:geometry:10}
Construct a rhombus whose diagonals measure $6$ cm and $4$ cm,
knowing that the diagonals of a rhombus cross at their middles at a
right angle: draw the diagonals \emph{first}, then join their four
ends.
\end{exercise}

\begin{exercise}[$\star\star$]\label{exo:g5:geometry:11}
A goat is tied to a stake $A$ by a $4$ m rope, in a flat field with
a straight fence passing $3$ m from $A$. Draw a plan (1 cm per
meter): which region can the goat reach? Where does the fence limit
it?
\end{exercise}
